% ---------------------------------------------------------------
% R3.07 SQL dans un langage de programmation -- preambule commun des diapositives
% Charte derivee de R5.A.08 Qualite de developpement (metropolis, dark tech)
% Palette IUT partagee (violet du logo + accent or)
% Utilisation :
%   \documentclass[aspectratio=169,10pt]{beamer}
%   % ---------------------------------------------------------------
% R3.07 SQL dans un langage de programmation -- preambule commun des diapositives
% Charte derivee de R5.A.08 Qualite de developpement (metropolis, dark tech)
% Palette IUT partagee (violet du logo + accent or)
% Utilisation :
%   \documentclass[aspectratio=169,10pt]{beamer}
%   % ---------------------------------------------------------------
% R3.07 SQL dans un langage de programmation -- preambule commun des diapositives
% Charte derivee de R5.A.08 Qualite de developpement (metropolis, dark tech)
% Palette IUT partagee (violet du logo + accent or)
% Utilisation :
%   \documentclass[aspectratio=169,10pt]{beamer}
%   % ---------------------------------------------------------------
% R3.07 SQL dans un langage de programmation -- preambule commun des diapositives
% Charte derivee de R5.A.08 Qualite de developpement (metropolis, dark tech)
% Palette IUT partagee (violet du logo + accent or)
% Utilisation :
%   \documentclass[aspectratio=169,10pt]{beamer}
%   \input{../../latex/common/r307-slides-preamble.tex}
%   \graphicspath{{../../latex/common/assets/}{img/}{figs/}}
% ---------------------------------------------------------------
\usetheme{metropolis}
\usepackage{appendixnumberbeamer}

\usepackage[T1]{fontenc}
\usepackage[utf8]{inputenc}
\usepackage[french]{babel}
\usepackage[scale=0.94]{tgheros}
\renewcommand{\familydefault}{\sfdefault}
\usepackage{amsmath,amssymb,amsfonts}
\usepackage{graphicx}
\usepackage{booktabs}
\usepackage{multirow}
\usepackage{tikz}
\usetikzlibrary{shapes.geometric, shapes, arrows, arrows.meta, positioning, calc, fit}
\usepackage{xcolor}
\usepackage{listings}
\usepackage[skins,listings]{tcolorbox}
\usepackage{qrcode}
\usepackage{hyperref}

% ---------------------------------------------------------------
% PALETTE (dark-mode tech: near-black canvas, violet (logo IUT) + gold accents)
% ---------------------------------------------------------------
\definecolor{qabg}{HTML}{0E1016}      % canvas background
\definecolor{qapanel}{HTML}{181B24}   % card / block panel
\definecolor{qapanelalt}{HTML}{20242F}% secondary panel tint
\definecolor{qaline}{HTML}{2B303D}    % hairline rules / table rules
\definecolor{qatext}{HTML}{ECEEF2}    % primary text (off-white)
\definecolor{qamuted}{HTML}{8B93A6}   % secondary / caption text
\definecolor{qaaccent}{HTML}{D93ED7}  % primary accent -- violet du logo IUT, eclairci pour fond sombre
\definecolor{qaaccent2}{HTML}{F2B84B} % secondary accent -- or/ambre
\definecolor{qacodekw}{HTML}{C792EA}  % code : mots-cles
\definecolor{qacodestr}{HTML}{C3E88D} % code : chaines
\definecolor{qacodecom}{HTML}{7F8AA3} % code : commentaires
\definecolor{qacodeann}{HTML}{F2B84B} % code : annotations / javadoc

\setbeamercolor{background canvas}{bg=qabg}
\setbeamercolor{normal text}{fg=qatext}
\setbeamercolor{palette primary}{bg=qabg, fg=qamuted}
\setbeamercolor{frametitle}{bg=qabg, fg=qatext}
\setbeamercolor{progress bar}{fg=qaaccent, bg=qaline}
\setbeamercolor{alerted text}{fg=qaaccent2}
\setbeamercolor{block title}{bg=qapanel, fg=qaaccent}
\setbeamercolor{block body}{bg=qapanel, fg=qatext}
\setbeamercolor{block title alerted}{bg=qapanel, fg=qaaccent2}
\setbeamercolor{block body alerted}{bg=qapanel, fg=qatext}
\setbeamercolor{block title example}{bg=qapanelalt, fg=qaaccent}
\setbeamercolor{block body example}{bg=qapanelalt, fg=qatext}
\setbeamercolor{example text}{fg=qaaccent}
\setbeamercolor{itemize item}{fg=qaaccent}
\setbeamercolor{itemize subitem}{fg=qaaccent2}
\setbeamercolor{enumerate item}{fg=qaaccent}
\setbeamercolor{enumerate subitem}{fg=qaaccent2}
\setbeamercolor{caption}{fg=qamuted}
\setbeamercolor{page number in head/foot}{fg=qamuted, bg=qabg}
\setbeamercolor{title}{fg=qatext}
\setbeamercolor{subtitle}{fg=qamuted}
\setbeamercolor{author}{fg=qatext}
\setbeamercolor{institute}{fg=qamuted}
\setbeamercolor{date}{fg=qamuted}
\setbeamercolor{section in toc}{fg=qatext}
\setbeamercolor{subsection in toc}{fg=qamuted}

\setbeamertemplate{blocks}[rounded][shadow=false]
\setbeamertemplate{itemize item}{\raisebox{0.15em}{\scriptsize\textcolor{qaaccent}{\ensuremath{\blacktriangleright}}}}
\setbeamertemplate{itemize subitem}{\raisebox{0.1em}{\tiny\textcolor{qaaccent2}{\ensuremath{\bullet}}}}

\setbeamerfont{frametitle}{series=\bfseries, size=\Large}
\setbeamerfont{title}{series=\bfseries, size=\Large}
\setbeamerfont{subtitle}{size=\normalsize}
\setbeamerfont{institute}{size=\scriptsize}
\setbeamerfont{date}{size=\scriptsize}
\setbeamerfont{author}{size=\small}

\metroset{sectionpage=none, subsectionpage=none, numbering=fraction, progressbar=frametitle}
\setbeamertemplate{section in toc}[sections numbered]

% ---------------------------------------------------------------
% LISTINGS -- theme sombre (lstlisting reste utilisable tel quel)
% \lstcode : police + couleur de base, a reutiliser dans les
%            surcharges locales (ex. basicstyle=\lstcode\scriptsize)
% ---------------------------------------------------------------
\newcommand{\lstcode}{\ttfamily\color{qatext}}

\lstdefinelanguage{TypeScript}{
  keywords={typeof, new, true, false, catch, function, return, null, switch, var,
    if, in, while, do, else, case, break, const, let, interface, class, extends,
    implements, enum, export, import, this, super, constructor, public, private,
    protected, static, async, await, throw, throws, try, finally, void},
  ndkeywords={boolean, any, number, string, int, long, double},
  sensitive=true,
  comment=[l]{//},
  morecomment=[s]{/*}{*/},
  morestring=[b]',
  morestring=[b]"
}

\lstset{
  language=Java,
  basicstyle=\lstcode\footnotesize,
  keywordstyle=\color{qacodekw}\bfseries,
  ndkeywordstyle=\color{qaaccent2}\bfseries,
  identifierstyle=\color{qatext},
  stringstyle=\color{qacodestr},
  commentstyle=\color{qacodecom}\itshape,
  morecomment=[s][\color{qacodeann}]{/**}{*/},
  backgroundcolor=\color{qapanelalt},
  rulecolor=\color{qaline},
  frame=single,
  framesep=3pt,
  xleftmargin=4pt,
  xrightmargin=4pt,
  keepspaces=true,
  columns=fullflexible,
  showstringspaces=false,
  breaklines=true,
  tabsize=2,
  escapeinside={(*@}{@*)},
}

% ---------------------------------------------------------------
% CARD-STYLE BOXES
% ---------------------------------------------------------------
\tcbset{
  figcard/.style={colback=white, colframe=white, boxrule=0pt, arc=4pt,
    left=3pt, right=3pt, top=3pt, bottom=3pt, boxsep=0pt,
    halign=center, valign=center},
  datacard/.style={colback=qapanel, colframe=qaline, boxrule=0.4pt, arc=3pt,
    left=8pt, right=8pt, top=6pt, bottom=6pt, coltext=qatext,
    before upper={\color{qatext}}, halign=center},
  codebox/.style={colback=qapanelalt, colframe=qaline, boxrule=0.4pt, arc=3pt,
    left=8pt, right=8pt, top=5pt, bottom=5pt, listing only,
    listing options={basicstyle=\ttfamily\small\color{qaaccent},
      breaklines=true, showstringspaces=false, columns=fullflexible,
      keepspaces=true, frame=none, backgroundcolor={}}}
}
\newtcolorbox{figcard}[1][]{figcard,#1}
\newtcolorbox{datacard}{datacard}
\newtcblisting{codebox}{codebox}

% ---------------------------------------------------------------
% HANDY MACROS
% ---------------------------------------------------------------
\newcommand{\hl}[1]{\alert{#1}}
\newcommand{\source}[1]{{\scriptsize\color{qamuted}#1}}

% Sommaire automatique en debut de section
\AtBeginSection[]{%
  \begin{frame}<beamer>{Plan du cours}
    \tableofcontents[currentsection]
  \end{frame}}

% Page de titre : \rtitleslide{numero de chapitre}{titre du chapitre}
% Les auteurs sont repris de \rauthors (modifiable avant \begin{document})
\newcommand{\rauthors}{Selena Lamari et Chouki Tibermacine}
\newcommand{\rtitleslide}[2]{%
\begin{frame}[plain,noframenumbering]
  \begin{center}
    \vspace{2mm}
    {\color{qaaccent}\rule{2.1cm}{1.2pt}}\\[6mm]
    {\Large\bfseries\color{qatext} R3.07}\\[1mm]
    {\Large\bfseries\color{qatext} SQL dans un langage de programmation}\\[5mm]
    {\large\color{qamuted} Chapitre #1}\\[1mm]
    {\large\color{qaaccent} #2}\\[6mm]
    \includegraphics[height=2.1cm]{iut_logo_transparent.png}

    \vspace{4mm}

    {\small\color{qatext} \rauthors}\\[1mm]
    {\scriptsize\color{qamuted} BUT 2 Informatique --- IUT de Vannes --- Universit\'e Bretagne Sud}
  \end{center}
\end{frame}
}

% ---- Diapo de fin « Questions ? » (même modèle que R5.A.08) ----
\newcommand{\rcontact}{Chouki.Tibermacine@univ-ubs.fr}
\newcommand{\rquestionslide}{%
\begin{frame}{Questions ?}
  \begin{center}
    \vspace{2mm}
    \begin{tikzpicture}
      \node[regular polygon, regular polygon sides=6, draw=qaline, thick,
        minimum size=4.6cm, rotate=90] (hexin) {};
      \node[regular polygon, regular polygon sides=6, draw=qaaccent, dashed,
        thick, minimum size=6.2cm, rotate=90] (hexout) {};
      \node[font=\fontsize{92}{92}\selectfont\bfseries, color=qaaccent2] at (hexin) {?};
    \end{tikzpicture}\\[5mm]
    {\small\color{qamuted} \rcontact}
  \end{center}
\end{frame}}

%   \graphicspath{{../../latex/common/assets/}{img/}{figs/}}
% ---------------------------------------------------------------
\usetheme{metropolis}
\usepackage{appendixnumberbeamer}

\usepackage[T1]{fontenc}
\usepackage[utf8]{inputenc}
\usepackage[french]{babel}
\usepackage[scale=0.94]{tgheros}
\renewcommand{\familydefault}{\sfdefault}
\usepackage{amsmath,amssymb,amsfonts}
\usepackage{graphicx}
\usepackage{booktabs}
\usepackage{multirow}
\usepackage{tikz}
\usetikzlibrary{shapes.geometric, shapes, arrows, arrows.meta, positioning, calc, fit}
\usepackage{xcolor}
\usepackage{listings}
\usepackage[skins,listings]{tcolorbox}
\usepackage{qrcode}
\usepackage{hyperref}

% ---------------------------------------------------------------
% PALETTE (dark-mode tech: near-black canvas, violet (logo IUT) + gold accents)
% ---------------------------------------------------------------
\definecolor{qabg}{HTML}{0E1016}      % canvas background
\definecolor{qapanel}{HTML}{181B24}   % card / block panel
\definecolor{qapanelalt}{HTML}{20242F}% secondary panel tint
\definecolor{qaline}{HTML}{2B303D}    % hairline rules / table rules
\definecolor{qatext}{HTML}{ECEEF2}    % primary text (off-white)
\definecolor{qamuted}{HTML}{8B93A6}   % secondary / caption text
\definecolor{qaaccent}{HTML}{D93ED7}  % primary accent -- violet du logo IUT, eclairci pour fond sombre
\definecolor{qaaccent2}{HTML}{F2B84B} % secondary accent -- or/ambre
\definecolor{qacodekw}{HTML}{C792EA}  % code : mots-cles
\definecolor{qacodestr}{HTML}{C3E88D} % code : chaines
\definecolor{qacodecom}{HTML}{7F8AA3} % code : commentaires
\definecolor{qacodeann}{HTML}{F2B84B} % code : annotations / javadoc

\setbeamercolor{background canvas}{bg=qabg}
\setbeamercolor{normal text}{fg=qatext}
\setbeamercolor{palette primary}{bg=qabg, fg=qamuted}
\setbeamercolor{frametitle}{bg=qabg, fg=qatext}
\setbeamercolor{progress bar}{fg=qaaccent, bg=qaline}
\setbeamercolor{alerted text}{fg=qaaccent2}
\setbeamercolor{block title}{bg=qapanel, fg=qaaccent}
\setbeamercolor{block body}{bg=qapanel, fg=qatext}
\setbeamercolor{block title alerted}{bg=qapanel, fg=qaaccent2}
\setbeamercolor{block body alerted}{bg=qapanel, fg=qatext}
\setbeamercolor{block title example}{bg=qapanelalt, fg=qaaccent}
\setbeamercolor{block body example}{bg=qapanelalt, fg=qatext}
\setbeamercolor{example text}{fg=qaaccent}
\setbeamercolor{itemize item}{fg=qaaccent}
\setbeamercolor{itemize subitem}{fg=qaaccent2}
\setbeamercolor{enumerate item}{fg=qaaccent}
\setbeamercolor{enumerate subitem}{fg=qaaccent2}
\setbeamercolor{caption}{fg=qamuted}
\setbeamercolor{page number in head/foot}{fg=qamuted, bg=qabg}
\setbeamercolor{title}{fg=qatext}
\setbeamercolor{subtitle}{fg=qamuted}
\setbeamercolor{author}{fg=qatext}
\setbeamercolor{institute}{fg=qamuted}
\setbeamercolor{date}{fg=qamuted}
\setbeamercolor{section in toc}{fg=qatext}
\setbeamercolor{subsection in toc}{fg=qamuted}

\setbeamertemplate{blocks}[rounded][shadow=false]
\setbeamertemplate{itemize item}{\raisebox{0.15em}{\scriptsize\textcolor{qaaccent}{\ensuremath{\blacktriangleright}}}}
\setbeamertemplate{itemize subitem}{\raisebox{0.1em}{\tiny\textcolor{qaaccent2}{\ensuremath{\bullet}}}}

\setbeamerfont{frametitle}{series=\bfseries, size=\Large}
\setbeamerfont{title}{series=\bfseries, size=\Large}
\setbeamerfont{subtitle}{size=\normalsize}
\setbeamerfont{institute}{size=\scriptsize}
\setbeamerfont{date}{size=\scriptsize}
\setbeamerfont{author}{size=\small}

\metroset{sectionpage=none, subsectionpage=none, numbering=fraction, progressbar=frametitle}
\setbeamertemplate{section in toc}[sections numbered]

% ---------------------------------------------------------------
% LISTINGS -- theme sombre (lstlisting reste utilisable tel quel)
% \lstcode : police + couleur de base, a reutiliser dans les
%            surcharges locales (ex. basicstyle=\lstcode\scriptsize)
% ---------------------------------------------------------------
\newcommand{\lstcode}{\ttfamily\color{qatext}}

\lstdefinelanguage{TypeScript}{
  keywords={typeof, new, true, false, catch, function, return, null, switch, var,
    if, in, while, do, else, case, break, const, let, interface, class, extends,
    implements, enum, export, import, this, super, constructor, public, private,
    protected, static, async, await, throw, throws, try, finally, void},
  ndkeywords={boolean, any, number, string, int, long, double},
  sensitive=true,
  comment=[l]{//},
  morecomment=[s]{/*}{*/},
  morestring=[b]',
  morestring=[b]"
}

\lstset{
  language=Java,
  basicstyle=\lstcode\footnotesize,
  keywordstyle=\color{qacodekw}\bfseries,
  ndkeywordstyle=\color{qaaccent2}\bfseries,
  identifierstyle=\color{qatext},
  stringstyle=\color{qacodestr},
  commentstyle=\color{qacodecom}\itshape,
  morecomment=[s][\color{qacodeann}]{/**}{*/},
  backgroundcolor=\color{qapanelalt},
  rulecolor=\color{qaline},
  frame=single,
  framesep=3pt,
  xleftmargin=4pt,
  xrightmargin=4pt,
  keepspaces=true,
  columns=fullflexible,
  showstringspaces=false,
  breaklines=true,
  tabsize=2,
  escapeinside={(*@}{@*)},
}

% ---------------------------------------------------------------
% CARD-STYLE BOXES
% ---------------------------------------------------------------
\tcbset{
  figcard/.style={colback=white, colframe=white, boxrule=0pt, arc=4pt,
    left=3pt, right=3pt, top=3pt, bottom=3pt, boxsep=0pt,
    halign=center, valign=center},
  datacard/.style={colback=qapanel, colframe=qaline, boxrule=0.4pt, arc=3pt,
    left=8pt, right=8pt, top=6pt, bottom=6pt, coltext=qatext,
    before upper={\color{qatext}}, halign=center},
  codebox/.style={colback=qapanelalt, colframe=qaline, boxrule=0.4pt, arc=3pt,
    left=8pt, right=8pt, top=5pt, bottom=5pt, listing only,
    listing options={basicstyle=\ttfamily\small\color{qaaccent},
      breaklines=true, showstringspaces=false, columns=fullflexible,
      keepspaces=true, frame=none, backgroundcolor={}}}
}
\newtcolorbox{figcard}[1][]{figcard,#1}
\newtcolorbox{datacard}{datacard}
\newtcblisting{codebox}{codebox}

% ---------------------------------------------------------------
% HANDY MACROS
% ---------------------------------------------------------------
\newcommand{\hl}[1]{\alert{#1}}
\newcommand{\source}[1]{{\scriptsize\color{qamuted}#1}}

% Sommaire automatique en debut de section
\AtBeginSection[]{%
  \begin{frame}<beamer>{Plan du cours}
    \tableofcontents[currentsection]
  \end{frame}}

% Page de titre : \rtitleslide{numero de chapitre}{titre du chapitre}
% Les auteurs sont repris de \rauthors (modifiable avant \begin{document})
\newcommand{\rauthors}{Selena Lamari et Chouki Tibermacine}
\newcommand{\rtitleslide}[2]{%
\begin{frame}[plain,noframenumbering]
  \begin{center}
    \vspace{2mm}
    {\color{qaaccent}\rule{2.1cm}{1.2pt}}\\[6mm]
    {\Large\bfseries\color{qatext} R3.07}\\[1mm]
    {\Large\bfseries\color{qatext} SQL dans un langage de programmation}\\[5mm]
    {\large\color{qamuted} Chapitre #1}\\[1mm]
    {\large\color{qaaccent} #2}\\[6mm]
    \includegraphics[height=2.1cm]{iut_logo_transparent.png}

    \vspace{4mm}

    {\small\color{qatext} \rauthors}\\[1mm]
    {\scriptsize\color{qamuted} BUT 2 Informatique --- IUT de Vannes --- Universit\'e Bretagne Sud}
  \end{center}
\end{frame}
}

% ---- Diapo de fin « Questions ? » (même modèle que R5.A.08) ----
\newcommand{\rcontact}{Chouki.Tibermacine@univ-ubs.fr}
\newcommand{\rquestionslide}{%
\begin{frame}{Questions ?}
  \begin{center}
    \vspace{2mm}
    \begin{tikzpicture}
      \node[regular polygon, regular polygon sides=6, draw=qaline, thick,
        minimum size=4.6cm, rotate=90] (hexin) {};
      \node[regular polygon, regular polygon sides=6, draw=qaaccent, dashed,
        thick, minimum size=6.2cm, rotate=90] (hexout) {};
      \node[font=\fontsize{92}{92}\selectfont\bfseries, color=qaaccent2] at (hexin) {?};
    \end{tikzpicture}\\[5mm]
    {\small\color{qamuted} \rcontact}
  \end{center}
\end{frame}}

%   \graphicspath{{../../latex/common/assets/}{img/}{figs/}}
% ---------------------------------------------------------------
\usetheme{metropolis}
\usepackage{appendixnumberbeamer}

\usepackage[T1]{fontenc}
\usepackage[utf8]{inputenc}
\usepackage[french]{babel}
\usepackage[scale=0.94]{tgheros}
\renewcommand{\familydefault}{\sfdefault}
\usepackage{amsmath,amssymb,amsfonts}
\usepackage{graphicx}
\usepackage{booktabs}
\usepackage{multirow}
\usepackage{tikz}
\usetikzlibrary{shapes.geometric, shapes, arrows, arrows.meta, positioning, calc, fit}
\usepackage{xcolor}
\usepackage{listings}
\usepackage[skins,listings]{tcolorbox}
\usepackage{qrcode}
\usepackage{hyperref}

% ---------------------------------------------------------------
% PALETTE (dark-mode tech: near-black canvas, violet (logo IUT) + gold accents)
% ---------------------------------------------------------------
\definecolor{qabg}{HTML}{0E1016}      % canvas background
\definecolor{qapanel}{HTML}{181B24}   % card / block panel
\definecolor{qapanelalt}{HTML}{20242F}% secondary panel tint
\definecolor{qaline}{HTML}{2B303D}    % hairline rules / table rules
\definecolor{qatext}{HTML}{ECEEF2}    % primary text (off-white)
\definecolor{qamuted}{HTML}{8B93A6}   % secondary / caption text
\definecolor{qaaccent}{HTML}{D93ED7}  % primary accent -- violet du logo IUT, eclairci pour fond sombre
\definecolor{qaaccent2}{HTML}{F2B84B} % secondary accent -- or/ambre
\definecolor{qacodekw}{HTML}{C792EA}  % code : mots-cles
\definecolor{qacodestr}{HTML}{C3E88D} % code : chaines
\definecolor{qacodecom}{HTML}{7F8AA3} % code : commentaires
\definecolor{qacodeann}{HTML}{F2B84B} % code : annotations / javadoc

\setbeamercolor{background canvas}{bg=qabg}
\setbeamercolor{normal text}{fg=qatext}
\setbeamercolor{palette primary}{bg=qabg, fg=qamuted}
\setbeamercolor{frametitle}{bg=qabg, fg=qatext}
\setbeamercolor{progress bar}{fg=qaaccent, bg=qaline}
\setbeamercolor{alerted text}{fg=qaaccent2}
\setbeamercolor{block title}{bg=qapanel, fg=qaaccent}
\setbeamercolor{block body}{bg=qapanel, fg=qatext}
\setbeamercolor{block title alerted}{bg=qapanel, fg=qaaccent2}
\setbeamercolor{block body alerted}{bg=qapanel, fg=qatext}
\setbeamercolor{block title example}{bg=qapanelalt, fg=qaaccent}
\setbeamercolor{block body example}{bg=qapanelalt, fg=qatext}
\setbeamercolor{example text}{fg=qaaccent}
\setbeamercolor{itemize item}{fg=qaaccent}
\setbeamercolor{itemize subitem}{fg=qaaccent2}
\setbeamercolor{enumerate item}{fg=qaaccent}
\setbeamercolor{enumerate subitem}{fg=qaaccent2}
\setbeamercolor{caption}{fg=qamuted}
\setbeamercolor{page number in head/foot}{fg=qamuted, bg=qabg}
\setbeamercolor{title}{fg=qatext}
\setbeamercolor{subtitle}{fg=qamuted}
\setbeamercolor{author}{fg=qatext}
\setbeamercolor{institute}{fg=qamuted}
\setbeamercolor{date}{fg=qamuted}
\setbeamercolor{section in toc}{fg=qatext}
\setbeamercolor{subsection in toc}{fg=qamuted}

\setbeamertemplate{blocks}[rounded][shadow=false]
\setbeamertemplate{itemize item}{\raisebox{0.15em}{\scriptsize\textcolor{qaaccent}{\ensuremath{\blacktriangleright}}}}
\setbeamertemplate{itemize subitem}{\raisebox{0.1em}{\tiny\textcolor{qaaccent2}{\ensuremath{\bullet}}}}

\setbeamerfont{frametitle}{series=\bfseries, size=\Large}
\setbeamerfont{title}{series=\bfseries, size=\Large}
\setbeamerfont{subtitle}{size=\normalsize}
\setbeamerfont{institute}{size=\scriptsize}
\setbeamerfont{date}{size=\scriptsize}
\setbeamerfont{author}{size=\small}

\metroset{sectionpage=none, subsectionpage=none, numbering=fraction, progressbar=frametitle}
\setbeamertemplate{section in toc}[sections numbered]

% ---------------------------------------------------------------
% LISTINGS -- theme sombre (lstlisting reste utilisable tel quel)
% \lstcode : police + couleur de base, a reutiliser dans les
%            surcharges locales (ex. basicstyle=\lstcode\scriptsize)
% ---------------------------------------------------------------
\newcommand{\lstcode}{\ttfamily\color{qatext}}

\lstdefinelanguage{TypeScript}{
  keywords={typeof, new, true, false, catch, function, return, null, switch, var,
    if, in, while, do, else, case, break, const, let, interface, class, extends,
    implements, enum, export, import, this, super, constructor, public, private,
    protected, static, async, await, throw, throws, try, finally, void},
  ndkeywords={boolean, any, number, string, int, long, double},
  sensitive=true,
  comment=[l]{//},
  morecomment=[s]{/*}{*/},
  morestring=[b]',
  morestring=[b]"
}

\lstset{
  language=Java,
  basicstyle=\lstcode\footnotesize,
  keywordstyle=\color{qacodekw}\bfseries,
  ndkeywordstyle=\color{qaaccent2}\bfseries,
  identifierstyle=\color{qatext},
  stringstyle=\color{qacodestr},
  commentstyle=\color{qacodecom}\itshape,
  morecomment=[s][\color{qacodeann}]{/**}{*/},
  backgroundcolor=\color{qapanelalt},
  rulecolor=\color{qaline},
  frame=single,
  framesep=3pt,
  xleftmargin=4pt,
  xrightmargin=4pt,
  keepspaces=true,
  columns=fullflexible,
  showstringspaces=false,
  breaklines=true,
  tabsize=2,
  escapeinside={(*@}{@*)},
}

% ---------------------------------------------------------------
% CARD-STYLE BOXES
% ---------------------------------------------------------------
\tcbset{
  figcard/.style={colback=white, colframe=white, boxrule=0pt, arc=4pt,
    left=3pt, right=3pt, top=3pt, bottom=3pt, boxsep=0pt,
    halign=center, valign=center},
  datacard/.style={colback=qapanel, colframe=qaline, boxrule=0.4pt, arc=3pt,
    left=8pt, right=8pt, top=6pt, bottom=6pt, coltext=qatext,
    before upper={\color{qatext}}, halign=center},
  codebox/.style={colback=qapanelalt, colframe=qaline, boxrule=0.4pt, arc=3pt,
    left=8pt, right=8pt, top=5pt, bottom=5pt, listing only,
    listing options={basicstyle=\ttfamily\small\color{qaaccent},
      breaklines=true, showstringspaces=false, columns=fullflexible,
      keepspaces=true, frame=none, backgroundcolor={}}}
}
\newtcolorbox{figcard}[1][]{figcard,#1}
\newtcolorbox{datacard}{datacard}
\newtcblisting{codebox}{codebox}

% ---------------------------------------------------------------
% HANDY MACROS
% ---------------------------------------------------------------
\newcommand{\hl}[1]{\alert{#1}}
\newcommand{\source}[1]{{\scriptsize\color{qamuted}#1}}

% Sommaire automatique en debut de section
\AtBeginSection[]{%
  \begin{frame}<beamer>{Plan du cours}
    \tableofcontents[currentsection]
  \end{frame}}

% Page de titre : \rtitleslide{numero de chapitre}{titre du chapitre}
% Les auteurs sont repris de \rauthors (modifiable avant \begin{document})
\newcommand{\rauthors}{Selena Lamari et Chouki Tibermacine}
\newcommand{\rtitleslide}[2]{%
\begin{frame}[plain,noframenumbering]
  \begin{center}
    \vspace{2mm}
    {\color{qaaccent}\rule{2.1cm}{1.2pt}}\\[6mm]
    {\Large\bfseries\color{qatext} R3.07}\\[1mm]
    {\Large\bfseries\color{qatext} SQL dans un langage de programmation}\\[5mm]
    {\large\color{qamuted} Chapitre #1}\\[1mm]
    {\large\color{qaaccent} #2}\\[6mm]
    \includegraphics[height=2.1cm]{iut_logo_transparent.png}

    \vspace{4mm}

    {\small\color{qatext} \rauthors}\\[1mm]
    {\scriptsize\color{qamuted} BUT 2 Informatique --- IUT de Vannes --- Universit\'e Bretagne Sud}
  \end{center}
\end{frame}
}

% ---- Diapo de fin « Questions ? » (même modèle que R5.A.08) ----
\newcommand{\rcontact}{Chouki.Tibermacine@univ-ubs.fr}
\newcommand{\rquestionslide}{%
\begin{frame}{Questions ?}
  \begin{center}
    \vspace{2mm}
    \begin{tikzpicture}
      \node[regular polygon, regular polygon sides=6, draw=qaline, thick,
        minimum size=4.6cm, rotate=90] (hexin) {};
      \node[regular polygon, regular polygon sides=6, draw=qaaccent, dashed,
        thick, minimum size=6.2cm, rotate=90] (hexout) {};
      \node[font=\fontsize{92}{92}\selectfont\bfseries, color=qaaccent2] at (hexin) {?};
    \end{tikzpicture}\\[5mm]
    {\small\color{qamuted} \rcontact}
  \end{center}
\end{frame}}

%   \graphicspath{{../../latex/common/assets/}{img/}{figs/}}
% ---------------------------------------------------------------
\usetheme{metropolis}
\usepackage{appendixnumberbeamer}

\usepackage[T1]{fontenc}
\usepackage[utf8]{inputenc}
\usepackage[french]{babel}
\usepackage[scale=0.94]{tgheros}
\renewcommand{\familydefault}{\sfdefault}
\usepackage{amsmath,amssymb,amsfonts}
\usepackage{graphicx}
\usepackage{booktabs}
\usepackage{multirow}
\usepackage{tikz}
\usetikzlibrary{shapes.geometric, shapes, arrows, arrows.meta, positioning, calc, fit}
\usepackage{xcolor}
\usepackage{listings}
\usepackage[skins,listings]{tcolorbox}
\usepackage{qrcode}
\usepackage{hyperref}

% ---------------------------------------------------------------
% PALETTE (dark-mode tech: near-black canvas, violet (logo IUT) + gold accents)
% ---------------------------------------------------------------
\definecolor{qabg}{HTML}{0E1016}      % canvas background
\definecolor{qapanel}{HTML}{181B24}   % card / block panel
\definecolor{qapanelalt}{HTML}{20242F}% secondary panel tint
\definecolor{qaline}{HTML}{2B303D}    % hairline rules / table rules
\definecolor{qatext}{HTML}{ECEEF2}    % primary text (off-white)
\definecolor{qamuted}{HTML}{8B93A6}   % secondary / caption text
\definecolor{qaaccent}{HTML}{D93ED7}  % primary accent -- violet du logo IUT, eclairci pour fond sombre
\definecolor{qaaccent2}{HTML}{F2B84B} % secondary accent -- or/ambre
\definecolor{qacodekw}{HTML}{C792EA}  % code : mots-cles
\definecolor{qacodestr}{HTML}{C3E88D} % code : chaines
\definecolor{qacodecom}{HTML}{7F8AA3} % code : commentaires
\definecolor{qacodeann}{HTML}{F2B84B} % code : annotations / javadoc

\setbeamercolor{background canvas}{bg=qabg}
\setbeamercolor{normal text}{fg=qatext}
\setbeamercolor{palette primary}{bg=qabg, fg=qamuted}
\setbeamercolor{frametitle}{bg=qabg, fg=qatext}
\setbeamercolor{progress bar}{fg=qaaccent, bg=qaline}
\setbeamercolor{alerted text}{fg=qaaccent2}
\setbeamercolor{block title}{bg=qapanel, fg=qaaccent}
\setbeamercolor{block body}{bg=qapanel, fg=qatext}
\setbeamercolor{block title alerted}{bg=qapanel, fg=qaaccent2}
\setbeamercolor{block body alerted}{bg=qapanel, fg=qatext}
\setbeamercolor{block title example}{bg=qapanelalt, fg=qaaccent}
\setbeamercolor{block body example}{bg=qapanelalt, fg=qatext}
\setbeamercolor{example text}{fg=qaaccent}
\setbeamercolor{itemize item}{fg=qaaccent}
\setbeamercolor{itemize subitem}{fg=qaaccent2}
\setbeamercolor{enumerate item}{fg=qaaccent}
\setbeamercolor{enumerate subitem}{fg=qaaccent2}
\setbeamercolor{caption}{fg=qamuted}
\setbeamercolor{page number in head/foot}{fg=qamuted, bg=qabg}
\setbeamercolor{title}{fg=qatext}
\setbeamercolor{subtitle}{fg=qamuted}
\setbeamercolor{author}{fg=qatext}
\setbeamercolor{institute}{fg=qamuted}
\setbeamercolor{date}{fg=qamuted}
\setbeamercolor{section in toc}{fg=qatext}
\setbeamercolor{subsection in toc}{fg=qamuted}

\setbeamertemplate{blocks}[rounded][shadow=false]
\setbeamertemplate{itemize item}{\raisebox{0.15em}{\scriptsize\textcolor{qaaccent}{\ensuremath{\blacktriangleright}}}}
\setbeamertemplate{itemize subitem}{\raisebox{0.1em}{\tiny\textcolor{qaaccent2}{\ensuremath{\bullet}}}}

\setbeamerfont{frametitle}{series=\bfseries, size=\Large}
\setbeamerfont{title}{series=\bfseries, size=\Large}
\setbeamerfont{subtitle}{size=\normalsize}
\setbeamerfont{institute}{size=\scriptsize}
\setbeamerfont{date}{size=\scriptsize}
\setbeamerfont{author}{size=\small}

\metroset{sectionpage=none, subsectionpage=none, numbering=fraction, progressbar=frametitle}
\setbeamertemplate{section in toc}[sections numbered]

% ---------------------------------------------------------------
% LISTINGS -- theme sombre (lstlisting reste utilisable tel quel)
% \lstcode : police + couleur de base, a reutiliser dans les
%            surcharges locales (ex. basicstyle=\lstcode\scriptsize)
% ---------------------------------------------------------------
\newcommand{\lstcode}{\ttfamily\color{qatext}}

\lstdefinelanguage{TypeScript}{
  keywords={typeof, new, true, false, catch, function, return, null, switch, var,
    if, in, while, do, else, case, break, const, let, interface, class, extends,
    implements, enum, export, import, this, super, constructor, public, private,
    protected, static, async, await, throw, throws, try, finally, void},
  ndkeywords={boolean, any, number, string, int, long, double},
  sensitive=true,
  comment=[l]{//},
  morecomment=[s]{/*}{*/},
  morestring=[b]',
  morestring=[b]"
}

\lstset{
  language=Java,
  basicstyle=\lstcode\footnotesize,
  keywordstyle=\color{qacodekw}\bfseries,
  ndkeywordstyle=\color{qaaccent2}\bfseries,
  identifierstyle=\color{qatext},
  stringstyle=\color{qacodestr},
  commentstyle=\color{qacodecom}\itshape,
  morecomment=[s][\color{qacodeann}]{/**}{*/},
  backgroundcolor=\color{qapanelalt},
  rulecolor=\color{qaline},
  frame=single,
  framesep=3pt,
  xleftmargin=4pt,
  xrightmargin=4pt,
  keepspaces=true,
  columns=fullflexible,
  showstringspaces=false,
  breaklines=true,
  tabsize=2,
  escapeinside={(*@}{@*)},
}

% ---------------------------------------------------------------
% CARD-STYLE BOXES
% ---------------------------------------------------------------
\tcbset{
  figcard/.style={colback=white, colframe=white, boxrule=0pt, arc=4pt,
    left=3pt, right=3pt, top=3pt, bottom=3pt, boxsep=0pt,
    halign=center, valign=center},
  datacard/.style={colback=qapanel, colframe=qaline, boxrule=0.4pt, arc=3pt,
    left=8pt, right=8pt, top=6pt, bottom=6pt, coltext=qatext,
    before upper={\color{qatext}}, halign=center},
  codebox/.style={colback=qapanelalt, colframe=qaline, boxrule=0.4pt, arc=3pt,
    left=8pt, right=8pt, top=5pt, bottom=5pt, listing only,
    listing options={basicstyle=\ttfamily\small\color{qaaccent},
      breaklines=true, showstringspaces=false, columns=fullflexible,
      keepspaces=true, frame=none, backgroundcolor={}}}
}
\newtcolorbox{figcard}[1][]{figcard,#1}
\newtcolorbox{datacard}{datacard}
\newtcblisting{codebox}{codebox}

% ---------------------------------------------------------------
% HANDY MACROS
% ---------------------------------------------------------------
\newcommand{\hl}[1]{\alert{#1}}
\newcommand{\source}[1]{{\scriptsize\color{qamuted}#1}}

% Sommaire automatique en debut de section
\AtBeginSection[]{%
  \begin{frame}<beamer>{Plan du cours}
    \tableofcontents[currentsection]
  \end{frame}}

% Page de titre : \rtitleslide{numero de chapitre}{titre du chapitre}
% Les auteurs sont repris de \rauthors (modifiable avant \begin{document})
\newcommand{\rauthors}{Selena Lamari et Chouki Tibermacine}
\newcommand{\rtitleslide}[2]{%
\begin{frame}[plain,noframenumbering]
  \begin{center}
    \vspace{2mm}
    {\color{qaaccent}\rule{2.1cm}{1.2pt}}\\[6mm]
    {\Large\bfseries\color{qatext} R3.07}\\[1mm]
    {\Large\bfseries\color{qatext} SQL dans un langage de programmation}\\[5mm]
    {\large\color{qamuted} Chapitre #1}\\[1mm]
    {\large\color{qaaccent} #2}\\[6mm]
    \includegraphics[height=2.1cm]{iut_logo_transparent.png}

    \vspace{4mm}

    {\small\color{qatext} \rauthors}\\[1mm]
    {\scriptsize\color{qamuted} BUT 2 Informatique --- IUT de Vannes --- Universit\'e Bretagne Sud}
  \end{center}
\end{frame}
}

% ---- Diapo de fin « Questions ? » (même modèle que R5.A.08) ----
\newcommand{\rcontact}{Chouki.Tibermacine@univ-ubs.fr}
\newcommand{\rquestionslide}{%
\begin{frame}{Questions ?}
  \begin{center}
    \vspace{2mm}
    \begin{tikzpicture}
      \node[regular polygon, regular polygon sides=6, draw=qaline, thick,
        minimum size=4.6cm, rotate=90] (hexin) {};
      \node[regular polygon, regular polygon sides=6, draw=qaaccent, dashed,
        thick, minimum size=6.2cm, rotate=90] (hexout) {};
      \node[font=\fontsize{92}{92}\selectfont\bfseries, color=qaaccent2] at (hexin) {?};
    \end{tikzpicture}\\[5mm]
    {\small\color{qamuted} \rcontact}
  \end{center}
\end{frame}}
